\documentclass[11pt]{article}
\usepackage[margin=25mm]{geometry}
\usepackage{fontspec}
\setmainfont{Latin Modern Roman}
\usepackage{amsmath,amssymb,mathtools}
\usepackage{unicode-math}
\setmathfont{Latin Modern Math}
\usepackage{xurl,hyperref}
\hypersetup{colorlinks=true,urlcolor=blue}
\setcounter{secnumdepth}{0}
\setlength{\parindent}{0pt}
\setlength{\parskip}{5pt}
\setlength{\emergencystretch}{3em}
\newcommand{\sref}[2]{\hyperlink{source:#1:#2}{[S#2]}}
\newcommand{\cataloguescope}[1]{\paragraph{Recorded target.}#1}
\begin{document}
\section{Zarankiewicz Problem}\label{statement}

\subsection{Definitions and mathematical statement}

\paragraph{Shared notation.}$\mathbb N=\{1,2,\ldots\}$, and $\mathbb Z,\mathbb Q,\mathbb R,\mathbb C$ denote the integers, rationals, reals and complex numbers. Any other conventions are specified locally.


Let $m,n,s,t\in\mathbb N$. A bipartite graph $G=(L\sqcup R,E)$ has fixed partite sets with $|L|=m$, $|R|=n$, and edges only between $L$ and $R$. In the standard partite-respecting Zarankiewicz convention, a forbidden $K_{s,t}$ consists of $S\subseteq L$, $T\subseteq R$, $|S|=s$, $|T|=t$, such that every pair in $S\times T$ is an edge. Equivalently, the $m\times n$ bipartite adjacency matrix is a zero--one matrix with no all-one $s\times t$ submatrix obtained by choosing rows and columns; the rows and columns need not be consecutive. Define
\[
 z(m,n;s,t)=\max\{|E(G)|:G\text{ has these partite sizes and no such }K_{s,t}\}.
\]
This fixes the roles of the two partite sets, as in the source's matrix definition; it gives $z(m,n;s,t)=z(n,m;t,s)$.
\paragraph{Statement recorded in the supplied source.}
Determine $z(m,n;s,t)$ for all positive integers $m,n,s,t$, including the unresolved exact finite values and sharp extremal behavior. The request is for the full four-parameter extremal function, rather than only the equal-part case or one fixed forbidden graph. \sref{statement}{2}
\subsection{Short English statement}
Determine the greatest number of edges between two prescribed vertex classes when no chosen s vertices in the first class and t in the second are all joined to each other.
\subsection{Sources}
\begin{description}
\item[\hypertarget{source:statement:1}{S1}] ULAM.AI and the UnsolvedMath Contributors, UnsolvedMath: A Curated Collection of Open Mathematics Problems, catalogue entry. CC BY 4.0. \url{https://huggingface.co/datasets/ulamai/UnsolvedMath}.
\item[\hypertarget{source:statement:2}{S2}] Vladimir Nikiforov, A Contribution to the Zarankiewicz Problem, arXiv:0903.5350v1 (2009), abstract and section “The matrix Zarankiewicz problem”, pages 1--2. \url{https://arxiv.org/pdf/0903.5350v1}.
\end{description}
\label{end:statement}
\end{document}
